\addcontentsline{toc}{section}{1.1　等加速度運動の一般解}

\begin{problemBox}{問題 1.1\quad 等加速度運動の一般解\hspace{1.2em}{\small\mdseries 〔基本〕}}
質量 $m$ の無人探査機が，宇宙空間を一直線上に運動している。軌道に沿って $x$ 軸をとり，探査機は $x$ 軸の正方向に一定の推進力 $F_0$ を受けるものとする。

時刻 $t=0$ において，探査機は原点 $x=0$ を速度 $v_0>0$ で通過した。探査機を質点とみなし，以下の問いに答えよ。

\begin{enumerate}
  \item 探査機の位置を $x(t)$，速度を
  $v(t)=\dfrac{dx}{dt}$ とするとき，運動方程式を示せ。

  \item (1) で得られた運動方程式を，速度 $v(t)$ に関する1階微分方程式とみなす。初期条件 $v(0)=v_0$ のもとでこの微分方程式を解き，速度 $v(t)$ を時間の関数として求めよ。

  \item (2) の結果と $v(t)=\dfrac{dx}{dt}$ の関係を用いて，位置 $x(t)$ に関する微分方程式を立てよ。これを初期条件 $x(0)=0$ のもとで解き，位置 $x(t)$ を時間の関数として求めよ。

  \item 探査機が原点から距離 $L>0$ だけ進んだときの速度 $v_L$ を，$m$，$F_0$，$v_0$，$L$ を用いて表せ。
\end{enumerate}
\end{problemBox}

\solutionheading
\begin{solutionparts}

\solutionpart{(1)}
運動方程式より，質量と加速度の積は物体に働く外力に等しくなります。加速度は $\dfrac{d^2x}{dt^2}$ または $\dfrac{dv}{dt}$ と表されます。

探査機に働く外力は $x$ 軸正の向きに一定の力 $F_0$ であるため，求める運動方程式は以下のようになります。
\begin{equation*}
  m\frac{d^2x}{dt^2}=F_0
\end{equation*}
あるいは，速度 $v$ を用いて以下のように表すこともできます。
\begin{equation*}
  m\frac{dv}{dt}=F_0
\end{equation*}

\solutionpart{(2)}
(1) で得られた運動方程式の両辺を $m>0$ で割ると，次のようになります。
\begin{equation*}
  \frac{dv}{dt}=\frac{F_0}{m}
\end{equation*}
この方程式は，両辺を時間 $t$ で積分することにより解くことができます。
\begin{equation*}
  v(t)=\int\frac{F_0}{m}\,dt=\frac{F_0}{m}t+C_1
\end{equation*}
ここで $C_1$ は積分定数です。初期条件 $v(0)=v_0$ を代入すると，
\begin{equation*}
  v(0)=\frac{F_0}{m}\cdot 0+C_1=v_0
\end{equation*}
となり，$C_1=v_0$ が得られます。したがって，速度 $v(t)$ の時間依存性は以下のように求まります。
\begin{equation*}
  \boxed{v(t)=v_0+\frac{F_0}{m}t}
\end{equation*}
加速度は一定値 $a=\dfrac{F_0}{m}$ です。

\solutionpart{(3)}
速度と位置の関係式 $v(t)=\dfrac{dx}{dt}$ に (2) の結果を代入すると，次の1階微分方程式が得られます。
\begin{equation*}
  \frac{dx}{dt}=v_0+\frac{F_0}{m}t
\end{equation*}
この方程式の両辺を時間 $t$ で積分します。
\begin{align*}
  x(t)
  &=\int\left(v_0+\frac{F_0}{m}t\right)dt \\
  &=v_0t+\frac{F_0}{2m}t^2+C_2
\end{align*}
ここで $C_2$ は積分定数です。初期条件 $x(0)=0$ を代入すると，
\begin{equation*}
  x(0)=v_0\cdot 0+\frac{F_0}{2m}\cdot 0^2+C_2=0
\end{equation*}
となり，$C_2=0$ が得られます。したがって，位置 $x(t)$ の時間依存性は以下のようになります。
\begin{equation*}
  \boxed{x(t)=v_0t+\frac{F_0}{2m}t^2}
\end{equation*}
これは一定加速度 $a=\dfrac{F_0}{m}$ の等加速度運動を表します。

\solutionpart{(4)}
探査機が距離 $L$ だけ進んだときの時刻を $t_L$ とします。このとき，$x(t_L)=L$ が成り立ちます。(3) の結果より，次の関係式が得られます。
\begin{equation*}
  v_0t_L+\frac{F_0}{2m}t_L^2=L
\end{equation*}
速度と位置の式から時刻 $t$ を消去します。(2) より，時刻 $t$ における速度 $v(t)$ の式から $t$ を孤立させます。
\begin{equation*}
  t=\frac{m}{F_0}\bigl(v(t)-v_0\bigr)
\end{equation*}
この式を (3) の位置の式に代入します。
\begin{align*}
  x(t)
  &=v_0\left\{\frac{m}{F_0}\bigl(v(t)-v_0\bigr)\right\}
    +\frac{F_0}{2m}
    \left\{\frac{m}{F_0}\bigl(v(t)-v_0\bigr)\right\}^2 \\
  &=\frac{mv_0}{F_0}\bigl(v(t)-v_0\bigr)
    +\frac{m}{2F_0}\bigl(v(t)-v_0\bigr)^2 \\
  &=\frac{m}{2F_0}\bigl(v(t)-v_0\bigr)
    \left\{2v_0+\bigl(v(t)-v_0\bigr)\right\} \\
  &=\frac{m}{2F_0}\bigl(v(t)-v_0\bigr)\bigl(v(t)+v_0\bigr) \\
  &=\frac{m}{2F_0}\bigl(v(t)^2-v_0^2\bigr).
\end{align*}
この式を整理すると，以下の関係式が得られます。
\begin{equation*}
  v(t)^2-v_0^2=\frac{2F_0}{m}x(t)
\end{equation*}
探査機が位置 $x(t_L)=L$ に達したときの速度を $v_L=v(t_L)$ とすると，
\begin{equation*}
  v_L^2=v_0^2+\frac{2F_0L}{m}
\end{equation*}
探査機は正の向きの力を受け続けて加速するため，$v_L>0$ です。したがって，平方根の正の値を採用することで，求める速度 $v_L$ は以下のように得られます。
\begin{equation*}
  \boxed{v_L=\sqrt{v_0^2+\frac{2F_0L}{m}}}
\end{equation*}

\end{solutionparts}
